\chapter{System, Data, and Methods}
\label{chap:methods}
\sloppy

Chapter~\ref{chap:theory} specified the relations that the campaign seeks to
observe and the claims it cannot support. This chapter translates that
framework into the executed persona assignments, network conditions,
propagation protocol, auditor regimes, and analysis rules. It first separates
inherited components from the thesis-specific work.

\section{Prior-work boundary}
The co-authored arXiv preprint introduced the auditor--node architecture, the Misinformation Index (MI), and the Misinformation Propagation Rate (MPR) \parencite{maurya2025simulating}. This thesis uses that system as an inherited instrument. Its thesis-specific layer comprises the climate corpus, the Pfeffer mapping, the Phase~2 factorial design, the composition analysis, and the D-net stress test. The verified allocation of individual contributions remains an explicit front-matter submission requirement.

\section{Simulation design}
\label{sec:methods-simulation}

This study used the repository's agent-based information-diffusion engine to
compare the propagation of climate and geoengineering claims across network
structures and persona compositions.  The design adapts the homogeneous and
heterogeneous conditions of the prior auditor--node preprint
\parencite{maurya2025simulating}; it does not reproduce that study's crime-news
corpus.  The present campaign instead uses climate-specific articles and belief
profiles.  All values reported in this section are taken from the executed
Phase~2 configurations, engine source, campaign plan, hand-off record, and raw
run schemas.

\subsection{Experimental units and factorial design}
\label{subsec:simulation-grid}

The principal experimental unit was one configuration--article cell.  The
factorial grid crossed eight network topologies with either 12 homogeneous
persona conditions or six heterogeneous persona mixes, two scoring modes, and
six seed articles:
\begin{equation}
  \underbrace{8 \times 12 \times 2 \times 6}_{\text{homogeneous}}
  +
  \underbrace{8 \times 6 \times 2 \times 6}_{\text{heterogeneous}}
  = 1{,}728 \text{ cells}.
  \label{eq:phase2-grid}
\end{equation}
The parser retained one analytical run per configuration (\(N=1\)); 77
experiment names had multiple archived directories. It selected the most
recently modified complete directory with at least one successful call, a rule
that is not equivalent to one execution per configuration. Accordingly, article cells
within a configuration are repeated observations from one simulation run, not
independent campaign replicates.  The design supports descriptive comparison
of the executed cells but not seed-to-seed uncertainty estimation.

The two scoring modes were continuous-only (\texttt{T2c}) and dual
(\texttt{T2d}).  In a dual run, the discrete score is the headline score and a
continuous score is stored as a sidecar for the same event.  The two modes are
separate measurement instruments and were not averaged.  Phase~1
discrete-only results were likewise not inserted into the Phase~2 grid.

\subsection{Articles and ground truth}
\label{subsec:simulation-articles}

The merged article library contains 24 records, but every cell in the main grid
uses the same six campaign articles: \texttt{scopex\_2017},
\texttt{chemtrails\_gates\_2018\_2021},
\texttt{sai\_geoengineering}, \texttt{paris\_agreement},
\texttt{climate\_consensus}, and \texttt{polar\_bears}.  The set spans
geoengineering research, a chemtrails conspiracy narrative, and established
climate-science claims.  Each article record contains an identifier, title,
source text, five yes/no audit questions, the corresponding Boolean
ground-truth answers, and source notes.  These structured question--answer
pairs provide the reference against which propagated text is scored.  The
remaining 18 records in the merged library are not additional observations in
the \(1{,}728\)-cell grid. \Cref{tab:article-provenance} records the scholarly
or institutional sources used to author the six campaign summaries and
Boolean keys. The executable stimulus is the checked-in JSON; the table does
not claim that the summaries are verbatim reprints of those sources.

\begin{table}[tbp]
  \centering
  \caption{Provenance of the six campaign articles. Each record has five
  yes/no questions and Boolean expected answers in
  \path{articles/merged.json}. Access date: 2026-09-19.}
  \label{tab:article-provenance}
  \small
  \begin{tabularx}{\textwidth}{@{}lX@{}}
    \toprule
    Identifier & Principal sources \\
    \midrule
    \texttt{scopex\_2017} &
      Keutsch Group SCoPEx documentation \parencite{keutsch2024scopex};
      Debnath et al.\ on the 2017 volume peak
      \parencite{debnath2023conspiracy}. \\
    \texttt{chemtrails\_gates\_2018\_2021} &
      Debnath et al.\ on chemtrails spillover
      \parencite{debnath2023conspiracy}. \\
    \texttt{sai\_geoengineering} &
      IPCC WG~I assessment of solar radiation modification
      \parencite{ipcc2021ar6}. \\
    \texttt{paris\_agreement} &
      UNFCCC Paris Agreement text \parencite{unfccc2015paris}. \\
    \texttt{climate\_consensus} &
      IPCC WG~I \parencite{ipcc2021ar6}; Cook et al.\ consensus analysis
      \parencite{cook2013consensus}. \\
    \texttt{polar\_bears} &
      IPCC WG~I Arctic/sea-ice assessment \parencite{ipcc2021ar6}. \\
    \bottomrule
  \end{tabularx}
\end{table}

Articles were processed sequentially within a configuration. Before processing
each article, the simulation cleared the article's prior history and all node
inboxes.
The complete source text was then inserted into the designated seed node at
tick zero.  No competing article group or intervention was enabled in the
Phase~2 configurations considered here.

\subsection{Persona conditions}
\label{subsec:simulation-personas}

The merged persona library contains 17 belief profiles.  A profile provides at
least an identifier, a display name, and a system prompt; the richer campaign
records also contain tags, emotional tone, ideological bias, topic expertise,
Debnath-type annotations, stylistic attributes, hashtag repertoires, and
provenance.  These profiles are theory-informed reductions of discourse types,
not HDBSCAN centroids reconstructed from the full Debnath Twitter corpus
\parencite{debnath2023conspiracy}.  Their purpose is controlled role
conditioning, not demographic representation.

The homogeneous arm (\(H\)) uses 12 profiles, one profile per configuration
and the same profile at all eight nodes.  They comprise four conspiracy
profiles (\texttt{conspiracy\_believer},
\texttt{conspiracy\_haarp\_weather},
\texttt{conspiracy\_depopulation}, and
\texttt{conspiracy\_climate\_piggyback}); three climate-action profiles
(\texttt{climate\_action\_advocate}, \texttt{climate\_justice\_youth}, and
\texttt{mitigation\_first\_policy}); and five science/environment profiles
(\texttt{environmental\_concern},
\texttt{ozone\_stratosphere\_specialist},
\texttt{biodiversity\_food\_security}, \texttt{climate\_scientist}, and
\texttt{science\_journalist}).

The heterogeneous arm (\(He\)) declares six ordered mixes of eight
identifiers. The sequential generators realise these orders directly.
In particular, \texttt{mix\_00} contains four conspiracy and four
climate/environment profiles, \texttt{mix\_01} contains two conspiracy and six
climate/environment profiles, and \texttt{mix\_02} contains no conspiracy
profile.  Mixes \texttt{mix\_03}--\texttt{mix\_05} introduce the additional
moderator, SCoPEx climate-action, and peripheral-skywatcher profiles recorded
in the merged library. Node assignment follows the order in each mix for the
linear, ring, ER, small-world, and scale-free generators. Echo-chamber,
polarised, and hierarchical configurations instead use cluster assignment.
For five of six mixes this filtering and cycling realises only four or five
unique personas, repeated across eight seats; \texttt{mix\_02} is the only
declared order realised as listed. Polarised persona assignment alternates by
node index and therefore does not align with the graph's two contiguous
halves. Consequently, \(H\) versus \(He\) is a bundled assignment contrast in
persona prevalence, diversity, placement, and sometimes assignment procedure,
not a controlled test of diversity.

\subsection{Network topologies}
\label{subsec:simulation-topologies}

Every main-grid graph contains eight directed nodes.  The eight topology
families and their executed parameters are summarised in
\Cref{tab:simulation-topologies}.  Edge trust is stored on the source node.
For generators that use persona homophily, identical persona identifiers keep
the unadjusted base trust. Otherwise the source code sets
\(\mathrm{clip}(\mathrm{baseTrust}-0.2+0.5\cdot\mathrm{Jaccard}(\mathrm{tags}),0.05,0.95)\).

\begin{table}[tbp]
  \centering
  \caption{Network topologies in the Phase~2 main grid.}
  \label{tab:simulation-topologies}
  \small
  \begin{tabularx}{\textwidth}{@{}lX@{}}
    \toprule
    Topology & Construction and configured parameters \\
    \midrule
    Linear chain &
      Directed path \(v_i\rightarrow v_{i+1}\), initial edge trust \(0.7\). \\
    Ring &
      Directed cycle \(v_i\rightarrow v_{(i+1)\bmod 8}\), initial edge trust
      \(0.7\). \\
    Random ER &
      Directed Erdős--Rényi graph with edge probability \(p=0.42\).
      The configuration records \texttt{minSeedOutDegree}=2 and a seeded RNG,
      but \texttt{buildRandomER} ignores both and samples unseeded
      \texttt{Math.random}; edge trust is \(0.3+U(0,0.5)\)
      \parencite{erdos1960evolution}. \\
    Small world &
      Watts--Strogatz-style directed ring lattice with \(k=4\) and rewiring
      probability \(\beta=0.1\) \parencite{watts1998collective}. \\
    Scale free &
      Preferential attachment with \(m=2\), beginning from a fully connected
      seed of \(m+1\) nodes \parencite{barabasi1999emergence}. \\
    Echo chamber &
      Two alternating chambers; within-/between-chamber edge probabilities
      \(0.75/0.08\), base trusts \(0.85/0.15\), and minimum seed out-degree 2. \\
    Polarised &
      Two halves; within-/between-group edge probabilities \(0.75/0.05\),
      base trusts \(0.88/0.10\), and minimum seed out-degree 2. \\
    Hierarchical &
      Bidirectional rooted 3-ary tree with downward/upward base trusts
      \(0.85/0.30\). \\
    \bottomrule
  \end{tabularx}
\end{table}

The linear chain is retained as the direct structural link to the earlier
prior linear-chain design.  The other generators test cyclic, random, clustered,
hub-dominated, ideologically partitioned, and authority--audience structures.
The eight network conditions in \Cref{tab:simulation-topologies} are experimental
mechanisms. Echo-chamber, polarised, and hierarchical generators are bespoke
treatments; their names do not imply validated platform models. Linear-chain
and ring configurations also use a more rewrite-heavy action policy and a
lower trust threshold than the other six generators, so chain- or
ring-versus-graph differences estimate a bundled protocol contrast.

\subsection{Seeding and propagation protocol}
\label{subsec:simulation-propagation}

Each article is seeded once into \texttt{node\_0} at tick zero. The seed message has
source \texttt{ORIGIN}, hop count zero, tick zero, the unmodified article text,
and an empty provenance chain.  Nodes are always active and have an inbox cap
of four messages per article.  Processing then proceeds synchronously for at
most eight ticks.  Messages produced during tick \(t\) are delivered only
after all nodes have completed that tick and are therefore available from the
next tick.  The simulation stops early when a tick produces no outgoing
message.

For a received message \(x\), the node first rejects the message if its hop
count has reached the configured limit or if effective source trust is below
the node's threshold. The receiver looks up its reverse relation to the source
and otherwise uses the hard-coded default \(0.5\); this need not equal trust
on the edge that carried the message. It then samples one of three configured actions:
\[
  A(x)\in\{\text{forward},\text{reinterpret},\text{drop}\}.
\]
Linear-chain and ring configurations use
\begin{equation}
  P(A)=\bigl(0.10,\;0.85,\;0.05\bigr),
  \qquad \tau=0.08,
  \label{eq:chain-actions}
\end{equation}
where the tuple is ordered as forward, reinterpret, and drop, and \(\tau\) is
the trust threshold. The other six main-grid configurations use
\begin{equation}
  P(A)=\bigl(0.35,\;0.50,\;0.15\bigr),
  \qquad \tau=0.15.
  \label{eq:graph-actions}
\end{equation}
A forward action preserves the received content.  A reinterpret action sends
the content to the node's persona-conditioned language model with an
instruction to rewrite it in 100--200 words and return only the rewritten
text.  A drop action records no outgoing text.  Forwarded and reinterpreted
messages are broadcast to all of the node's outgoing neighbours with hop count
incremented by one.
The four-message inbox cap silently discards excess deliveries; capacity loss
is not logged as a sampled drop.

\Cref{fig:simulation-pseudocode} summarises the executed protocol.
\begin{figure}[tbp]
  \centering
  \begin{minipage}{0.94\textwidth}
    \small
    \textbf{Input:} one configuration and the six fixed seed articles.
    \begin{enumerate}
      \item For each article \(a\), clear inboxes and prior history for \(a\),
      then deliver \(a\) to \texttt{node\_0} with \(h=0,t=0\).
      \item For \(t=1,\ldots,8\), process every message at every active node.
      If \(h\geq8\) or source trust is below \(\tau_v\), record
      \textsc{drop}.  Otherwise sample
      \(A\in\{\textsc{forward},\textsc{reinterpret},\textsc{drop}\}\).
      For \textsc{reinterpret}, rewrite the message using the node persona and
      \texttt{gpt-4o-mini}.  For either non-drop action, queue the resulting
      text for all out-neighbours with \(h\leftarrow h+1\).
      \item Deliver all queued messages synchronously.  Stop early if the
      queue is empty; otherwise continue with the next tick.
      \item After propagation of all six articles, score every forward or
      reinterpret output against its article's five ground-truth items.
    \end{enumerate}
  \end{minipage}
  \caption{Synchronous propagation of six articles under an eight-tick budget,
  followed by post-hoc auditing. Messages delivered after tick eight are not
  processed. The figure shows the executed protocol, not a validated social
  process.}
  \label{fig:simulation-pseudocode}
\end{figure}

The configuration includes \texttt{relationEvolution=true} and
\texttt{trustDelta}=0.05.  In the executed source, however, propagation of all
six articles is completed before the audit phase begins.  The auditor then
raises or lowers stored source trust by \(0.05\) according to whether the
event's misinformation index is at most or above 3.  Thus these post-hoc trust
updates cannot alter the already-completed propagation trajectories in the
same run.  Belief-state updates, frame analysis, strategic-agent behaviour,
network evolution, opinion dynamics, institutional trust, competitive groups,
and interventions are disabled in the raw Phase~2 metadata.

\subsection{Agent and auditor models}
\label{subsec:simulation-model}

Both rewriting agents and the auditor use the mutable alias
\texttt{gpt-4o-mini}. Chat Completions requests used temperature \(0.7\),
\texttt{max\_tokens}=700, a 120-second socket timeout, no API seed, no response
schema, and no retry or backoff. The runner normally used concurrency two.
Response identifiers, model fingerprints, and raw auditor responses were not
archived. The agent is conditioned by the assigned persona's system prompt.
The auditor receives the
propagated text, each of the five article questions, and the expected yes/no
answers.  Scores are therefore LLM-as-judge measurements rather than human
annotations \parencite{zheng2023judge}.  Auditing occurs after propagation and
only for events whose action is forward or reinterpret and whose output text
is present.

For discrete scoring, question \(j\) receives
\(d_j\in\{-1,0,1\}\), denoting contradiction/distortion, missing information,
or a correct statement, respectively.  For \(m=5\) questions, the discrete
misinformation index is
\begin{equation}
  \mathrm{MI}_{\mathrm{disc}}
  = \sum_{j=1}^{m}\mathbb{I}(d_j\leq 0)
  = n_{\mathrm{missing}}+n_{\mathrm{incorrect}},
  \qquad \mathrm{MI}_{\mathrm{disc}}\in\{0,\ldots,5\}.
  \label{eq:mm-mi-discrete}
\end{equation}
For continuous scoring, the auditor returns
\(c_j\in[0,1]\), where 1 denotes fully correct, 0.5 missing or neutral
information, and 0 contradiction.  The implemented score is
\begin{equation}
  \overline{c}=\frac{1}{m}\sum_{j=1}^{m}c_j,
  \qquad
  \mathrm{MI}_{\mathrm{cont}}=m(1-\overline{c}).
  \label{eq:mm-mi-continuous}
\end{equation}
The two nominal 0--5 ranges are not commensurate. A missing discrete item adds
1, whereas a neutral or missing continuous score of 0.5 adds 0.5. Five
missing items therefore yield 5 discretely but 2.5 continuously. The common
\(\mathrm{MI}>3\) cut has different meanings under the two instruments.

In dual mode, the discrete and continuous auditor calls are issued in parallel
for the same text.  The event's top-level \texttt{misinfoIndex} is
\(\mathrm{MI}_{\mathrm{disc}}\).  The raw event additionally stores both score
objects, the absolute diagnostic gap
\begin{equation}
  \Delta_{\mathrm{dual}}
  =\left|\mathrm{MI}_{\mathrm{disc}}-\mathrm{MI}_{\mathrm{cont}}\right|,
  \label{eq:dual-gap}
\end{equation}
and Pearson agreement after mapping the discrete item scores from
\(\{-1,0,1\}\) to \(\{0,0.5,1\}\).  The gap is a measurement diagnostic, not a
third misinformation scale.

\subsection{Eight-tick horizon and computational budget}
\label{subsec:eight-hop-budget}

The prior linear-chain study used 30 positions or hops
\parencite{maurya2025simulating}. Phase~2 limits each graph to eight nodes,
eight processing ticks, and a configured maximum of eight message hops. This is an explicit
computational-cost reduction: propagation can require an LLM rewrite for each
reinterpretation, followed by one auditor call per scored event in continuous
mode and two in dual mode.  Dense and cyclic graphs can generate many events
per tick, so retaining the 30-step horizon while expanding to eight
topologies, 288 configurations, six articles, and dual auditing would multiply
paid model calls. The eight-tick horizon therefore prioritises factorial
coverage over long-tail cascade depth. Messages delivered after tick eight
are not processed; the archived event records therefore contain audited input
hops 0--7 and no hop-8 events.

This truncation changes the estimand.  It measures early diffusion and
semantic change within a short, common budget; it does not establish that a
cascade has converged, and it weakens any claim of late-stage
propagation.  The number eight is also unrelated to Debnath et al.'s
eight-word skip-gram context window \parencite{debnath2023conspiracy}.  The
latter is a word-embedding hyperparameter, not a social-network hop protocol.

\subsection{Raw records and derived cell summaries}
\label{subsec:simulation-schema}

Each timestamped run directory contains \texttt{metadata.json},
\texttt{state.json}, \texttt{graph\_topology.json}, one
\texttt{results\_<articleId>.json} file per article, a
\texttt{human\_eval\_template.csv}, and one state file per node.  The raw
schemas were used as follows:
\begin{description}
  \item[\texttt{metadata.json}] run status, fully merged configuration,
  model-usage counts, and per-article result pointers;
  \item[\texttt{state.json}] propagation and audit checkpoints, completed
  article lists, current tick, errors, and last-update time;
  \item[\texttt{graph\_topology.json}] realised node labels, directed edges,
  and initial edge trusts;
  \item[\texttt{nodes/node\_*.json}] node and persona identifiers, model,
  relations, node parameters, inbox, action counters, and event history;
  \item[\texttt{results\_<articleId>.json}] node summaries and aggregate
  network, provenance, frame, and IFD metrics.
\end{description}
An event history record contains the tick, article, source node, hop count,
action, input and output text, misinformation index, reason, provenance,
timestamp, and---when auditing succeeds---the IFD score object.  This
event-level record, rather than the aggregate result file, is the parser's
source for Phase~2 headline scores.

For article \(a\), let \(E_a^\star\) be the set of history events with a
non-null top-level misinformation index.  The parsed cell headline is the
event-weighted mean
\begin{equation}
  \overline{\mathrm{MI}}_a
  = \frac{1}{|E_a^\star|}
    \sum_{e\in E_a^\star}\mathrm{MI}_e.
  \label{eq:cell-mean-mi}
\end{equation}
The parser also reports a mean of node-level means.  It is retained as a
separate quantity because it weights nodes equally rather than weighting all
events equally.  A cell with at most one scored event is marked dead and
hatched. Configuration-level \texttt{llmCalls>0} does not prove that the
particular article cell received an auditor call. Dead cells are retained in
the ledger, excluded from live means, and never imputed as zero.

\subsection{API failures and scoring coverage}
\label{subsec:api-failures}

Across the 288 selected directories, \texttt{llmUsage.calls} records 203,928
successful responses and \texttt{llmUsage.errors} records 64,270 parseable
HTTP responses with status at least 400. These counters are not retries or
duplicate reruns: the client has no retry loop and each selected directory is
counted once. They are also not a request ledger, because request identifiers
were not retained and concurrent dual calls can hide a sibling rejection.
At least one HTTP error affected 121 runs.

Rewrite-request failure caused the engine to propagate the input text
unchanged while retaining action \texttt{reinterpret}. Auditor-request failure
left MI null. Of 129,664 audit-eligible final-state events, 93,470 were scored
and 36,194 remained null; null events were excluded from means and
\(k^{*}\), and could change hatch status. The parser's all-correct fallback
for malformed auditor output is a latent fail-open risk, but the selected log
sections contained no parse-fallback warnings and node files contained no
recoverable scored-payload anomalies. The observed fallback count is therefore
zero, not an estimated source of score bias.

Only five affected experiment names have a complete zero-error alternative
run. Stochastic reruns were not silently substituted. Canonical tables,
identity-output exclusion, unaffected-cell, and zero-error-run analyses are
reported separately; none is a causal correction.

\subsection{Reproducibility and limitations}
\label{subsec:simulation-reproducibility}

The campaign records \texttt{graphRandomSeed: 42}, all merged configurations,
the realised topology, node histories, checkpoints, model-usage totals, and
the complete parser inputs.  The graph seed is consumed by the echo-chamber
and polarised builders.  Exact replay must nevertheless not be claimed.  In
the checked source revision, the ER, small-world, and scale-free builders, as
well as action sampling, call JavaScript's unseeded \texttt{Math.random};
the ER builder does not consume the additional seeded-options argument passed
by the simulation wrapper.  Language-model generation and judging can also
vary across calls.  Together with \(N=1\), these facts make the archived
realised graphs and event logs the reproducible record of what was executed,
while a fresh run is a protocol replication rather than a bitwise
reproduction.

\subsection{Campaign isolation}
\label{subsec:simulation-isolation}

Phase~2 was isolated at both configuration and output level.  New
configurations reside in
\texttt{thesisExperiment/configs/phase2/}; raw runs in
\texttt{thesisExperiment/runs\_phase2/}; and parsed outputs in
\texttt{thesisExperiment/results\_phase2/}.  The Phase~1 directories
\texttt{thesisExperiment/runs/},
\texttt{thesisExperiment/results/tables/}, and
\texttt{thesisExperiment/configs/full/} were designated read-only and were not
overwritten.  Parser selection is by experiment name and retains the latest
complete timestamped run with positive LLM-call usage.  This separation
prevents the earlier discrete \(12\times12\) campaign, probes, and custom
Debnath-network runs from silently entering the main \(1{,}728\)-cell
factorial analysis.


\section{Measurement and analysis}
\label{sec:methods-measurement}

The preceding section defined the simulation and raw event records. This
section now defines the analytical units, estimands, exclusion rule, and
comparisons applied to those records.

\subsection{Analytical scope and unit of observation}
\label{sec:mm-scope}

The measurement analysis used four nested units: audit item, scored event,
node--article MPR, and configuration--article cell. All scores are
LLM-as-judge measurements produced by \texttt{gpt-4o-mini}; they are neither
human annotations nor observations from Twitter.

Four nested measurement units must be distinguished:
\begin{enumerate}
  \item an \emph{item} is one of the five ground-truth questions;
  \item a \emph{scored event} is one agent's audited handling of one article;
  \item a \emph{node MPR} is the mean event-level misinformation index for one node and
        article; and
  \item an \emph{article cell} is one configuration--article combination.
\end{enumerate}
The primary cell statistic is \texttt{meanMI}, the arithmetic mean of event-level MI
over scored events.  The parser also records \texttt{meanNodeMPR}, the mean of the
node-level means.  These quantities use the same event-level instrument but weight
nodes differently; they are retained as separate columns and are not substituted for
one another.

\subsection{IFD scoring and the MI--MPR hierarchy}
\label{sec:mm-ifd}

The engine labels the auditor record \emph{IFD}.  For each event, the record contains a
misinformation index (MI), correct, missing, and incorrect rates
(\(\mathrm{CR}\), \(\mathrm{MR}\), and \(\mathrm{IR}\)), and auxiliary fields.  The
analysis uses MI as the event score.  It does not report hop-wise CR, MR, or IR because
those fields are absent from the live analysis tables.

\subsubsection{Discrete instrument}
\label{sec:mm-discrete}

For question \(j\), the discrete auditor returns
\[
  s_j\in\{-1,0,+1\},
\]
where \(+1\) denotes a correct answer, \(0\) a missing answer, and \(-1\) an incorrect
answer.  With \(m=5\) questions, let \(n_C,n_M,n_I\) denote the corresponding counts.
The implementation computes
\[
  \mathrm{CR}=\frac{n_C}{m},\qquad
  \mathrm{MR}=\frac{n_M}{m},\qquad
  \mathrm{IR}=\frac{n_I}{m},
\]
and
\begin{equation}
  \mathrm{MI}_{d}=n_M+n_I.
  \label{eq:mi-discrete}
\end{equation}
Thus \(\mathrm{MI}_{d}\in\{0,1,2,3,4,5\}\).  Missing and incorrect answers contribute
equally to MI even though the IFD record preserves their separate rates.  The
dual-scoring campaign, denoted T2d, uses this discrete MI as its top-level
\texttt{misinfoIndex} and hence as its headline instrument.

\subsubsection{Continuous instrument}
\label{sec:mm-continuous}

The continuous auditor returns one score \(q_j\in[0,1]\) per question.  Its correctness
rate and MI are
\begin{equation}
  \mathrm{CR}_{c}=\frac{1}{m}\sum_{j=1}^{m}q_j,
  \qquad
  \mathrm{MI}_{c}=m(1-\mathrm{CR}_{c}).
  \label{eq:mi-continuous}
\end{equation}
For the five-question articles, \(\mathrm{MI}_{c}\) is a real-valued score on the same
nominal \(0\)--\(5\) range.  Scores at least \(0.67\) are counted as correct and scores
at most \(0.33\) as incorrect when the continuous IFD record derives CR/MR/IR counts;
intermediate scores are counted as missing.  The continuous-only campaign, denoted
T2c, makes one continuous auditor call per scored event and stores
\(\mathrm{MI}_{c}\) as the event's headline \texttt{misinfoIndex}.

\subsubsection{MPR aggregation}
\label{sec:mm-mpr}

For node \(v\), article \(a\), and instrument \(r\in\{d,c\}\), the node-level
misinformation propagation rate (MPR) is
\begin{equation}
  \mathrm{MPR}^{(r)}_{v,a}
  =\frac{1}{|E_{v,a}|}\sum_{e\in E_{v,a}}\mathrm{MI}^{(r)}_e,
  \label{eq:mm-node-mpr}
\end{equation}
where \(E_{v,a}\) is the set of scored events for that node and article.  At cell level,
\texttt{meanMI} instead averages over all scored events directly:
\begin{equation}
  \overline{\mathrm{MI}}^{(r)}_{\mathrm{cell}}
  =\frac{1}{|E_{\mathrm{cell}}|}
   \sum_{e\in E_{\mathrm{cell}}}\mathrm{MI}^{(r)}_e.
  \label{eq:cell-mi}
\end{equation}
The parser's \texttt{meanNodeMPR} is
\(\frac{1}{|V_{\mathrm{scored}}|}\sum_v\mathrm{MPR}_{v,a}\).  It gives each reached
node equal weight, whereas \texttt{meanMI} gives each scored event equal weight.
The chapter therefore uses \texttt{meanMI} as the cell headline and presents
\texttt{meanNodeMPR} only as a companion summary.

\subsection{Dual scoring and instrument separation}
\label{sec:mm-dual}

In a T2d event the auditor makes two separate LLM calls in parallel: one discrete
call and one continuous call.  The returned top-level IFD object mirrors the discrete
record.  A nested \texttt{ifd.dual} sidecar stores the discrete record, the continuous
record, an absolute gap, and an agreement coefficient:
\begin{align}
  g_e
    &=\left|\mathrm{MI}_{d,e}-\mathrm{MI}_{c,e}\right|,\\
  a_e
    &=\operatorname{cor}\!\left(
       \frac{s_{1,e}+1}{2},\ldots,\frac{s_{m,e}+1}{2};
       q_{1,e},\ldots,q_{m,e}\right).
\end{align}
The normalisation maps discrete item scores \(-1,0,+1\) to \(0,0.5,1\), after which
the code computes Pearson's correlation with the continuous item scores.  The stored
gap is \emph{absolute}; it is not a signed discrete-minus-continuous difference.

This design yields two campaign instruments and one within-event diagnostic:
\begin{itemize}
  \item T2c headline: continuous MI from continuous-only events;
  \item T2d headline: discrete MI from dual-run events; and
  \item T2d sidecar: continuous MI, absolute gap, and item-score agreement on those
        same dual-run events.
\end{itemize}
The sidecar continuous score is not the T2c headline because it was obtained in a
different campaign call on a different realised event.  Likewise, the absolute gap is
a scorer-disagreement diagnostic, not a third MPR.  Discrete and continuous means are
therefore shown in separate columns or panels, never averaged into a common headline.
Cross-instrument differences are explicitly labelled as instrument contrasts and
cannot identify a change in misinformation while holding the measurement process
fixed.

\subsection{\texorpdfstring{Network threshold \(k^{*}\)}{Network threshold k*}}
\label{sec:mm-kstar}

The parser forms a network mean MI at every observed tick \(t\):
\[
  \overline{\mathrm{MI}}_t
  =\frac{1}{|E_t|}\sum_{e\in E_t}\mathrm{MI}_e.
\]
It then defines
\begin{equation}
  k^{*}
  =\min\left\{t:
       \overline{\mathrm{MI}}_t>3
       \ \land\
       \overline{\mathrm{MI}}_u>3
       \text{ for every later observed }u
     \right\}.
  \label{eq:kstar}
\end{equation}
A temporary crossing that later returns to \(3\) or below does not qualify. If no
crossing sustained through the last later scored observation is observed,
\(k^{*}\) is undefined rather than zero. The observation window ends at eight
ticks, so a crossing at the final scored tick is right-censored and does not
establish persistence beyond the window. T2d \(k^{*}\) is
calculated from the discrete series; T2c \(k^{*}\) from the continuous series.  Their
rates remain separate.  This threshold measures simulated semantic degradation under
the study's MI rule; it is not a Twitter ignition time or an empirical firestorm
duration.

\subsection{Hatch policy and analysis population}
\label{sec:mm-hatch}

A cell is classified as dead when it contains at most one scored event:
\[
  \texttt{dead}\iff \texttt{nScored}\leq 1.
\]
If the run nevertheless records positive LLM usage, the parser labels the row
\texttt{hatched\_dead\_after\_llm}.  Such rows remain in the harvest and in the dead-cell
ledger, but are hatched in plots and excluded from live-cell means.  They are neither
deleted nor imputed as \(\mathrm{MI}=0\).  Consequently, a dead heterogeneous cell is
not evidence that heterogeneity prevented propagation.  Missing and in-progress rows,
if present, would likewise not be empirical zeros.

The official harvest contains 292 hatched-dead cells out of 1,728 (16.9\%): 96 in
T2c\_H, 98 in T2d\_H, 50 in T2c\_He, and 48 in T2d\_He. Of these, 232 have
zero scored events and 60 have one scored event. The 60 one-event cells
contain observed MI but are excluded from live means by the hatch rule. The
four live populations are
therefore 480, 478, 238, and 240 cells, respectively.  All reported means, medians,
paired-cell contrasts, and sign counts use the live subset unless a count is
explicitly described as a full-harvest or dead-cell count.

\subsection{Pre-specified comparison structure}
\label{sec:mm-comparisons}

The five comparison families in \Cref{tab:mm-c1-c5} were designed to make the held and varied dimensions
visible.  A ``same topology'' comparison means that summaries are aligned by one of
the eight topology labels; it does not imply that independently executed T2c and T2d
runs share identical event histories.

\begin{table}[htbp]
  \centering
  \caption{Comparison families and admissible interpretation.}
  \label{tab:mm-c1-c5}
  \small
  \begin{tabular}{p{0.08\textwidth}p{0.31\textwidth}p{0.24\textwidth}p{0.27\textwidth}}
    \hline
    ID & Compared slices & Held/varied dimensions & Interpretation rule\\
    \hline
    C1 &
    T2c\_H vs.\ T2c\_He; T2d\_H vs.\ T2d\_He &
    topology and instrument held; identity arrangement varied &
    H--He contrasts are estimated separately for each instrument.\\

    C2 &
    T2c\_H vs.\ T2d\_H &
    topology and homogeneous arm held; instrument varied &
    an instrument comparison, not a common-scale treatment effect.\\

    C3 &
    T2c\_He vs.\ T2d\_He &
    topology and heterogeneous arm held; instrument varied &
    same restriction as C2; dual sidecar is reported separately.\\

    C4 &
    T2c\_H vs.\ T2d\_He; T2d\_H vs.\ T2c\_He &
    topology held; identity arrangement and instrument both varied &
    two labelled, confounded contrasts; no pooled C4 headline.\\

    C5 &
    live cells concatenated within a labelled slice or composition
    subset; D-net reported separately &
    cell-level table union across generated topologies &
    descriptive pooled summaries only; no physical super-graph.\\
    \hline
  \end{tabular}
\end{table}

For C1, H and He means are first calculated within topology and mode.  Both
cell-pooled and topology-equal summaries are retained because unequal live counts can
change weighting.  C2 and C3 additionally use paired live cells joined on topology,
persona or mix, and article.  A pair enters only when both corresponding cells are
live; a dead cell on either side is never replaced by zero.  C4 deliberately crosses
identity arrangement and instrument.  Its two directions are shown separately because
averaging them would manufacture a third, uninterpretable MPR.

\subsection{Pooled table union versus a physical super-graph}
\label{sec:mm-pooling}

In C5, ``pooled'' means row concatenation.  Let
\(\mathcal{L}_{r,z,g}\) be the live article-cell rows for instrument \(r\), identity arm
\(z\), and topology \(g\).  The all-topology pool is the disjoint table union
\[
  \mathcal{P}_{r,z}
  =\biguplus_{g=1}^{8}\mathcal{L}_{r,z,g}.
\]
The pooled mean weights each live article cell equally.  The topology-equal companion
first calculates one mean per topology and then gives each of the eight topology means
weight \(1/8\).  Reporting both exposes weighting sensitivity.

No edges are added between runs, no node identifiers are merged, and no propagation is
rerun on the union.  It is therefore incorrect to interpret C5 as a graph with 64
nodes, as a union of adjacency matrices, or as a synthetic super-cascade.  The D-net
analysis is also not part of this pool.  D-net is a separately executed 63-node,
228-directed-edge custom topology based on a hashtag co-occurrence object.  The edge
container is not a hydrated retweet cascade, and the object has no empirical
Twitter MPR.

\subsection{Descriptive and exploratory statistical procedures}
\label{sec:mm-statistics}

The primary analysis is descriptive.  It reports cell counts, live and dead counts,
arithmetic means of \texttt{meanMI} and \texttt{meanNodeMPR}, medians of paired
differences where relevant, \(k^{*}\) occurrence rates among live cells, dead-cell
rates over all harvested cells, and absolute dual-gap and agreement means from the
dual sidecar.  Grouped summaries are produced by topology, article, persona or mix,
identity arm, and instrument.  Heat-map entries with no live cell are left missing,
not displayed as zero.

The analysis reports sign counts, medians, ranges, denominators, and rank
coefficients as descriptive diagnostics. It does not use \(p\)-values as
evidence because the eight network conditions are fixed design strata, pooled
cells are nested within runs, and no sampling distribution over independent
campaign replicates exists. No causal claim, confidence interval, or
population-level effect is inferred.

\subsection{D-net structural validation}
\label{sec:mm-dnet-validation}

\subsubsection{Comparable cascade-shape metrics}

The validation code compares three structural quantities:
\begin{description}
  \item[Depth] the longest breadth-first-search path from the inferred root;
  \item[Breadth] the maximum number of reached nodes at any depth level; and
  \item[Structural virality] the mean pairwise shortest-path distance among reached
        nodes after treating propagation edges as undirected.
\end{description}
Cascade size is reported for context but excluded from the structural-similarity
score.  Simulated propagation trees are reconstructed from event provenance for
forward and reinterpret actions.  The implementation also records simulation span in
ticks; the empirical hashtag object has no usable tweet clock, so ticks cannot be
converted to hours.

For structural metric \(j\in\{\text{depth},\text{breadth},\text{virality}\}\), the
relative error is
\[
  e_j=\left|\frac{x^{\mathrm{sim}}_j}{x^{\mathrm{real}}_j}-1\right|,
\]
and the aggregate structural similarity is
\begin{equation}
  S_{\mathrm{struct}}
  =\max\left(0,1-\frac{1}{3}\sum_j e_j\right).
  \label{eq:struct-sim}
\end{equation}
The implementation additionally marks a metric as a scalar match when
\(e_j<0.30\).  For D-net, the empirical values were depth 4, breadth 33, and structural
virality 2.7; the corresponding means over 16 simulated cascades were 2.5625,
40.9375, and 1.7963.  Only breadth met the 30\% rule, yielding
\(S_{\mathrm{struct}}=0.6885\).  These are cascade-shape comparisons, not comparisons
of simulated MI with an empirical Twitter MI.

\subsubsection{KS and Jensen--Shannon comparisons}

For each of depth, breadth, and structural virality, the validation module compares
the empirical and simulated value arrays with an approximate two-sample
Kolmogorov--Smirnov (KS) test.  The KS statistic is the maximum distance between the
two empirical cumulative distribution functions.  A small statistic indicates closer
empirical distributions; the code reports a Kolmogorov-distribution approximation to
the \(p\)-value.

The same arrays are discretised into 20 common-width bins for Jensen--Shannon (JS)
divergence.  With base-2 logarithms, JS lies in \([0,1]\): zero denotes identical
binned distributions and one complete separation.  The code labels a distributional
match only when KS \(p>0.05\) \emph{and} JS \(<0.10\).  It defines the aggregate
distributional score as
\begin{equation}
  D= \max\left(0,1-\frac{1}{3}
       \sum_j \mathrm{JS}_j\right).
  \label{eq:distribution-score}
\end{equation}

The D-net comparison has one empirical object and 16 simulated cascades.  The observed
KS statistics were 1.0000 for depth, 0.8125 for breadth, and 1.0000 for structural
virality, with approximate \(p\)-values 0.1104, 0.2954, and 0.1104.  JS divergence was
1.0000 for all three metrics, so no metric met the joint match rule and \(D=0\).
Because \(n_{\mathrm{real}}=1\), the KS \(p\)-values are severely underpowered and do
not establish distributional equivalence.  In particular, failure to reject is not
evidence that the empirical and simulated distributions match.

\subsubsection{Digital Twin Fidelity Score}

The code combines structural similarity, distributional score, and a bounded content
correlation into the Digital Twin Fidelity Score (DTFS):
\begin{equation}
  \mathrm{DTFS}
  =0.40S_{\mathrm{struct}}+0.40D+0.20\rho_{\mathrm{content}},
  \label{eq:dtfs}
\end{equation}
with each component clipped to \([0,1]\). These weights, component definitions,
and the \(0.70\) threshold are implementation choices from the repository, not
externally validated criteria. DTFS is reported only to document that the
implemented check failed. The implementation's validation threshold
is \(\mathrm{DTFS}\geq0.70\).  D-net has
\(S_{\mathrm{struct}}=0.6885\), \(D=0\), and
\(\rho_{\mathrm{content}}=0\), giving
\[
  \mathrm{DTFS}=0.2754<0.70.
\]
Accordingly, \texttt{isValidated} is \texttt{false}.  This failure concerns the
implemented structural/distributional digital-twin check.  It does not turn the
hashtag graph into a retweet cascade, and it does not license an empirical MPR claim:
Debnath's data provide no empirical \(0\)--\(5\) MI/MPR against which the auditor
scores could be validated.

\subsection{Source register and reproducibility boundary}
\label{sec:mm-sources}

All operational definitions and numerical counts in this chapter are traceable to the
following repository artefacts:
\begin{itemize}
  \item IFD formulas, dual calls, absolute gap, and agreement:
        \path{src/Auditor.js};
  \item event collection, cell means, node means, \(k^{*}\), and hatch policy:
        \path{thesisExperiment/scripts/parse_phase2.js};
  \item harvested counts and the \(N=1\) parse-selection rule:
        \path{thesisExperiment/results_phase2/summary.json};
  \item C1--C5 joins, pooled summaries, and descriptive procedures:
        \path{thesisExperiment/analysis_full/analyze_full.py} and
        \path{thesisExperiment/analysis_full/tables/};
  \item structural metric extraction:
        \path{src/ValidationMetrics.js};
  \item KS, JS, structural similarity, and DTFS definitions:
        \path{src/ValidationComparison.js}; and
  \item D-net validation inputs and outputs:
        \path{thesisExperiment/results_phase2/debnath_compare.json}.
\end{itemize}
Phase~1 tables are not joined to these Phase~2 headlines because they use a different
campaign and schema.  D-net rows are not concatenated into the 1,728-cell factorial
grid.  These boundaries preserve the separation among instruments, campaigns, and
graph objects throughout the analysis.

The Results therefore proceed in four steps: auditor-regime sensitivity
(RQ1), named persona compositions (RQ2), topology-conditioned same-regime
contrasts with dead-cell rates (RQ3), and the separate D-net stress test
(RQ4). All means are live-cell descriptions; dead-cell rates are reported as
a distinct outcome.
\fussy
